% ===========================================================================
% Wing + Wheel RANS - setup and verification report
%
% Build:  ./build.sh          (regenerates figures, then runs latexmk)
%         latexmk -pdf report  (text only, reuse existing figures)
%
% Every number comes from generated/case_data.tex, which is written by
% scripts/make_tex_data.py out of the case's results.json and logs. Nothing
% in here is typed by hand.
% ===========================================================================
\documentclass[11pt,a4paper]{article}

\usepackage[utf8]{inputenc}
\usepackage[T1]{fontenc}
\usepackage{lmodern}
\usepackage[margin=22mm]{geometry}
\usepackage{graphicx}
\usepackage{booktabs}
\usepackage{longtable}
\usepackage{array}
\usepackage{amsmath}
% siunitx is not in every TeX installation, and this report needs almost
% nothing from it. Minimal stand-ins, so no extra package is required.
\providecommand{\si}[1]{\ensuremath{\mathrm{#1}}}
\providecommand{\SI}[2]{\ensuremath{#1\,\mathrm{#2}}}
\providecommand{\sisetup}[1]{}
\usepackage{xcolor}
\usepackage{enumitem}
\usepackage{caption}
\usepackage[hidelinks]{hyperref}
\usepackage{fancyhdr}

\setlist[itemize]{topsep=2pt, itemsep=1pt, parsep=0pt, leftmargin=*}
\captionsetup{font=small, labelfont=bf, skip=4pt}
\graphicspath{{figures/}}

\pagestyle{fancy}
\fancyhf{}
\lhead{\small Wing + Wheel RANS, OpenFOAM v\ofApi}
\rhead{\small \caseName}
\cfoot{\small \thepage}
\renewcommand{\headrulewidth}{0.3pt}

\definecolor{warn}{RGB}{176,58,46}
\definecolor{okc}{RGB}{30,110,60}
\newcommand{\ok}[1]{\textcolor{okc}{#1}}
\newcommand{\bad}[1]{\textcolor{warn}{#1}}

% Include a figure only if it exists, so a missing image never breaks the
% build - it leaves a visible placeholder instead.
\newcommand{\figOrNote}[4]{%
  \begin{figure}[htbp]\centering
  \IfFileExists{figures/#1.png}{\includegraphics[width=#2\linewidth]{#1}}%
    {\fbox{\parbox{0.8\linewidth}{\centering\vspace{8mm}\textbf{%
      \bad{figure \texttt{\detokenize{#1}.png} not generated}}\vspace{8mm}}}}%
  \caption{#3}\label{fig:#4}
  \end{figure}}

\input{generated/case_data}

\title{\vspace{-12mm}\textbf{Front wing in ground effect beside a rotating
wheel}\\[2mm]
\large Steady RANS setup and verification report\\[1mm]
\normalsize OpenFOAM v\ofApi\ (native) \textbf{·} \texttt{snappyHexMesh} +
\texttt{simpleFoam} \textbf{·} case \caseName}
\author{}
\date{}

\begin{document}
\maketitle
\thispagestyle{fancy}
\vspace{-10mm}

% ---------------------------------------------------------------------------
\begin{center}
\fcolorbox{warn}{warn!6}{\parbox{0.95\linewidth}{\small
\textbf{\bad{State of this report.}}
The setup, the mesh and the whole figure and reporting chain are complete and
verified. \textbf{The solution is not.} This build stops at
\textbf{\nIter\ SIMPLE iterations} of the \nIterations\ configured, so every
flow quantity in it --- forces, $C_p$, $C_pT$, $y^+$, $Q$ --- is a
\emph{partially converged} field and must not be quoted as a result. The
residuals are still at $10^{-2}$ (section~7).

Mesh-side numbers (cell counts, layer coverage, \texttt{checkMesh}, timings)
are final and can be relied on.

To finish it: run the solve to completion and rebuild, see
section~\ref{sec:finish}.
}}
\end{center}
\vspace{2mm}

% ===========================================================================
\section{Case at a glance}

\begin{center}\small
\begin{tabular}{@{}llll@{}}
\toprule
\multicolumn{2}{@{}l}{\textbf{Configuration}} &
\multicolumn{2}{l}{\textbf{Mesh}}\\
\midrule
wing half-span      & $S/c = \spanC$     & cells            & \nCells\\
ride height         & $h/c = \heightC$   & background cells & \nBaseCells\\
angle of attack     & \aoaDeg$^\circ$ & layer coverage & \layerCoverage\,\%\\
wheel width         & $W/c = \widthC$    & \texttt{checkMesh} clean & \checkMeshOk\\
\midrule
\multicolumn{2}{@{}l}{\textbf{Flow}} &
\multicolumn{2}{l}{\textbf{Cost, \npRanks\ MPI ranks}}\\
\midrule
freestream  & \Uinf\,\si{m/s}                  & meshing   & \meshMinutes\,min\\
chord       & \chord\,\si{m}                   & solving   & \solveMinutes\,min\\
$Re_c$      & \ReC                             & per iteration & \secPerIter\,s\\
$\rho$      & \rhoInf\,\si{kg/m^3}             & iterations & \nIter\\
$\nu$       & \nuInf\,\si{m^2/s}               & converged & \converged\\
\bottomrule
\end{tabular}
\end{center}

\paragraph{Headline result.}
Force \emph{coefficient-areas}, mean over the last 50 iterations.
$A_{\mathrm{ref}} = \SI{1}{m^2}$ on purpose, so these are $C_L\!\cdot\!A$ and
$C_D\!\cdot\!A$ in \si{m^2} --- the quantity the Diasinos papers tabulate.
Negative lift is downforce.

\begin{center}\small
\begin{tabular}{@{}lrrr@{}}
\toprule
 & $C_L\!\cdot\!A$ [\si{m^2}] & $C_D\!\cdot\!A$ [\si{m^2}] &
 spread of $C_L\!\cdot\!A$\\
\midrule
wing  & \clWing  & \cdWing  & \spreadWing\,\%\\
wheel & \clWheel & \cdWheel & \spreadWheel\,\%\\
total & \clTotal & \cdTotal & \spreadTotal\,\%\\
\bottomrule
\end{tabular}
\end{center}

\noindent
Mean $y^+$: wing \yplusWing, wheel \yplusWheel, maximum \yplusMax.

% ===========================================================================
\section{Software and reproducibility}

\begin{itemize}
\item \textbf{OpenFOAM v\ofApi} (ESI/OpenCFD), built from source. Native
  only --- no commercial solver, no in-house libraries.
\item Mesh \texttt{snappyHexMesh}, solve \texttt{simpleFoam} (segregated,
  SIMPLEC), turbulence \texttt{kOmegaSST}.
\item Geometry generated parametrically from the Diasinos specification;
  420 configurations available.
\item Everything else is \texttt{python3} standard library. No
  \texttt{numpy}/\texttt{pandas} needed to run a case.
\end{itemize}

\paragraph{To reproduce this case.}
\begin{verbatim}
cp env/project.env.example env/project.env   # set OPENFOAM_BASHRC and NP
cd Diasinos_Geometry_Generator && .venv/bin/python -m geometry_generator.generate
cd ../RANS_Simulations
python3 prepare_case.py --span 1.42 --height 0.13 --aoa 8 --width 0.63
cd S1.42_h0.13_AOA8_W0.63 && ./Allrun
\end{verbatim}

% ===========================================================================
\section{Geometry}

\begin{itemize}
\item Coordinate system: $x$ streamwise (wheel centre at $x=0$), $y$ spanwise
  outboard, $z$ up with the ground at $z=0$. Half model, symmetry at $y=0$.
\item Chord $c=\SI{75}{mm}$. Inverted NACA 4412, blunt trailing edge
  \SI{0.5}{mm}. Wheel diameter $1.17c$, shoulder radius $0.067c$.
\item \textbf{The generator writes millimetres}; everything in the
  simulation is metres. \texttt{Allmesh} scales by 0.001 on the way in.
\item The \emph{plinth} is a collar that deliberately pierces the ground
  plane (down to $z=\SI{-6}{mm}$) so the contact patch is a clean cut
  instead of a tangential singularity.
\end{itemize}

\figOrNote{geom_overview}{0.92}{Geometry and the ground plane. The wing sits
in ground effect ahead of and inboard of the wheel; the plinth passes through
the ground.}{geomov}

\figOrNote{geom_front}{0.92}{Front view, looking downstream. The endplate
(light blue) falls inside the wheel silhouette: at $S/c=\spanC$ and
$W/c=\widthC$ the spanwise gap is $-0.45c$, i.e. wing and wheel
overlap.}{geomfront}

\figOrNote{geom_plan}{0.86}{Plan view. $y=0$ is the symmetry plane.}{geomplan}

\begin{figure}[htbp]\centering
\IfFileExists{figures/geom_wing_regions.png}
  {\includegraphics[width=0.49\linewidth]{geom_wing_regions}}{}%
\IfFileExists{figures/geom_wheel_regions.png}
  {\includegraphics[width=0.49\linewidth]{geom_wheel_regions}}{}%
\caption{Named STL regions. \texttt{snappyHexMeshDict} refers to these names
directly to set refinement level and patch name, so renaming a region in the
generator breaks the mesh setup.}\label{fig:regions}
\end{figure}

\paragraph{Known geometry limitations.}
\begin{itemize}
\item The wing is \textbf{prismatic}, extruded between two spanwise
  stations. Geometrically exact for a constant section; it does mean the
  triangles have a very high aspect ratio, which \texttt{surfaceCheck}
  reports as poor quality. Not an error.
\item \bad{Two duplicate triangles} on the endplate outer edge --- a defect
  in the generator. Two facets out of 1142; \texttt{snappyHexMesh} tolerates
  them. Reported by \texttt{Allmesh} and deliberately non-fatal.
\item Both STLs consist of two disjoint parts (wing + endplate, wheel +
  plinth), so \texttt{surfaceCheck} says ``not closed''. Expected: no
  \texttt{mode inside} refinement region is attached to either body.
\item Wheel tessellation was raised from 180 to 900 circumferential segments,
  so the facet is \SI{0.31}{mm} against the \SI{0.293}{mm} cell. At 180
  segments the facets were \SI{1.53}{mm} --- five times coarser than the cell,
  i.e. the mesh resolved the wheel as a 180-gon with a visibly wavy surface.
  \emph{This did not improve the layer coverage on the tread}
  (63.4\,\% before and after); it was done for geometric fidelity alone.
\end{itemize}

% ===========================================================================
\section{Domain and boundary conditions}

\figOrNote{geom_domain}{0.92}{Computational domain,
$\SI{1.894}{m}\times\SI{0.675}{m}\times\SI{0.638}{m}
= 25.3c\times 9c\times 8.5c$. Grey: moving ground. Yellow: symmetry
plane.}{domain}

\begin{center}\small
\begin{tabular}{@{}lllll@{}}
\toprule
patch / group & type & $U$ & $p$ & $k,\omega,\nu_t$\\
\midrule
\texttt{inlet}  & patch & \texttt{fixedValue} (10\,0\,0) & \texttt{zeroGradient} & turbulent inlet\\
\texttt{outlet} & patch & \texttt{pressureInletOutletVelocity} & \texttt{fixedValue} 0 & \texttt{inletOutlet}\\
\texttt{symmetry} & symmetry & \texttt{symmetry} & \texttt{symmetry} & \texttt{symmetry}\\
\texttt{side}, \texttt{sky} & patch & \texttt{slip} & \texttt{slip} & \texttt{slip} / \texttt{calculated}\\
\texttt{ground} & wall & \texttt{fixedValue} (10\,0\,0) & \texttt{zeroGradient} & wall functions\\
\texttt{wingGroup} & wall & \texttt{noSlip} & \texttt{zeroGradient} & wall functions\\
\texttt{wheelRotating} & wall & \texttt{rotatingWallVelocity} & \texttt{zeroGradient} & wall functions\\
\texttt{wheelGround} & wall & \texttt{fixedValue} (10\,0\,0) & \texttt{zeroGradient} & wall functions\\
\bottomrule
\end{tabular}
\end{center}

\begin{itemize}
\item \textbf{Rolling road}: the belt translates downstream at the
  freestream speed, so the ground is at rest relative to the far-field air.
\item \textbf{Rotating wheel}: $\omega = \SI{-227.92}{rad/s}$ about
  $(0\,1\,0)$ through $(0,0,\SI{0.043875}{m})$.
  \texttt{rotatingWallVelocity} imposes $\omega\times r$ without moving the
  mesh --- correct for steady RANS. The sign follows from requiring the
  contact point to move with the belt:
  $\omega_y = -U_\infty/R$. Reversing it spins the wheel backwards and still
  converges happily.
\item The \texttt{wheel-plinth} patch is in \texttt{wheelGround}, not
  \texttt{wheelRotating}: it is the collar through the ground and moves with
  the belt.
\item \texttt{side} and \texttt{sky} are \texttt{patch}, not \texttt{wall},
  on purpose --- a wall there would enter the \texttt{meshWave} wall-distance
  field that \texttt{kOmegaSST} blends on, and would appear in the force
  groups.
\end{itemize}

% ===========================================================================
\section{Mesh}

\begin{itemize}
\item Background: $101\times36\times34$ \textbf{cubic} cells of
  \SI{18.75}{mm} ($=0.25c$), \nBaseCells\ cells. Cubic matters: each
  refinement level halves the edge, so the level-to-size table only holds
  from a cube.
\item Level sizes: L4 \SI{1.17}{mm}, L5 \SI{0.586}{mm}, L6 \SI{0.293}{mm},
  L7 \SI{0.146}{mm}.
\item Per-region surface levels, distance refinement around both bodies
  (L5 within \SI{15}{mm}, L4 within \SI{30}{mm}), and nine
  \texttt{searchableBox} zones.
\item Final mesh \nCells\ cells.
\end{itemize}

\begin{center}\small
\begin{tabular}{@{}lrr@{}}
\toprule
refinement level & cells & share [\%]\\
\midrule
\leveltable
\bottomrule
\end{tabular}
\end{center}

\figOrNote{mesh_slice_y_overview}{0.95}{Mesh on a spanwise slice through the
wheel centre. The nested refinement zones and the wake box are
visible.}{meshy}

\figOrNote{mesh_slice_y_gap}{0.95}{Same slice, wing trailing edge to wheel
front --- the region the whole study is about.}{meshgap}

\figOrNote{mesh_slice_x_overview}{0.86}{Slice through the wheel axis
($x=0$). Looking downstream, so $y$ increases to the left.}{meshx}

\figOrNote{mesh_slice_z}{0.86}{Mesh \SI{1.2}{mm} above the belt. The
contact-patch refinement box is the fine square under the
wheel.}{meshz}

\subsection{Prism layers}

Four layers from \SI{0.14}{mm} with expansion 1.3 (total \SI{0.87}{mm}) on
wing and wheel, three from \SI{0.25}{mm} on the ground. This targets
$y^+\approx 2.5$, which is why the wall functions are the all-$y^+$ variants.

\figOrNote{mesh_layers_tread}{0.80}{Wheel tread at the top. Prism layers on
the \SI{0.31}{mm}-tessellated tread.}{laytread}

\figOrNote{mesh_layers_tread_front}{0.80}{Wheel tread, upstream
face.}{laytreadf}

\figOrNote{mesh_layers_shoulder}{0.80}{Wheel shoulder, inboard. The
\SI{5}{mm} shoulder radius is the tightest convex curvature on the
wheel.}{layshoulder}

\figOrNote{mesh_layers_contact}{0.80}{Contact patch. The plinth cuts the
ground plane at $z=0$; this is where the layers fight each other.}{laycontact}

\figOrNote{mesh_wing_te}{0.80}{Blunt trailing edge, \SI{0.5}{mm} thick. No
prism layers --- see below.}{laywingte}

\begin{center}\small
\begin{tabular}{@{}lrrrr@{}}
\toprule
patch & faces & layers wanted & layers built & thickness [\%]\\
\midrule
\layertable
\bottomrule
\end{tabular}
\end{center}

\paragraph{Reading the layer table.}
\begin{itemize}
\item Overall coverage \layerCoverage\,\%.
\item The patches in \textbf{bold} got \bad{no layers at all}. These are
  \texttt{wing-TE}, \texttt{wing-endplate\_TE} and \texttt{wheel-plinth} ---
  sharp convex edges and the plinth/ground intersection. Asking for fewer and
  thinner layers does not change it: the limit is topology, not thickness.
\item \textbf{This is accepted.} \texttt{snappyHexMesh} is known for
  incomplete layer coverage at sharp edges and where two layered surfaces
  meet. All three patches sit in L7 cells (\SI{0.146}{mm}), so the near-wall
  spacing is comparable to the layered regions anyway.
\item The number to judge the near-wall treatment by is $y^+$
  (section~\ref{sec:results}), not the layer count.
\end{itemize}

\subsection{Mesh quality}

Quality limits are deliberately relaxed (\texttt{maxNonOrtho} 75,
\texttt{maxInternalSkewness} 6, \texttt{minTetQuality} $-10^{30}$), carried
over from the previous HELYX setup. Enforcing a positive tet quality with
absolute layer thicknesses on curved surfaces costs most of the layer
coverage.

\texttt{checkMesh} clean: \textbf{\checkMeshOk} (\nCheckMeshFail\ findings).
\begin{itemize}
\checkmeshlist
\end{itemize}

\noindent
These are the expected consequence of the relaxed limits. The solver runs
stably on this mesh; if it ever diverges, \texttt{meshQualityDict} is the
first thing to question.

% ===========================================================================
\section{Physics and numerics}

\begin{center}\small
\begin{tabular}{@{}ll@{}}
\toprule
turbulence model & \texttt{kOmegaSST}\\
$\nu_t$ wall function & \texttt{nutUSpaldingWallFunction} (continuous all $y^+$)\\
$k$ wall function & \texttt{kLowReWallFunction}\\
$\omega$ wall function & \texttt{omegaWallFunction}\\
\midrule
momentum convection & \texttt{bounded Gauss linearUpwindV grad(U)}\\
turbulence convection & \texttt{bounded Gauss upwind} (first order, deliberate)\\
gradients & \texttt{cellLimited Gauss linear 1}\\
laplacian & \texttt{Gauss linear corrected}\\
wall distance & \texttt{meshWave}\\
\midrule
pressure--velocity & SIMPLEC (\texttt{consistent yes})\\
$p$ solver & GAMG, \texttt{tolerance} $10^{-7}$, \texttt{relTol} 0.01\\
$U,k,\omega$ & \texttt{smoothSolver} GaussSeidel\\
relaxation & $U$ 0.9, $k$ and $\omega$ 0.7\\
\texttt{residualControl} & $p$, $U$: $10^{-5}$; $k$, $\omega$: $10^{-4}$\\
initialisation & \texttt{potentialFoam}, divergence-free start\\
\bottomrule
\end{tabular}
\end{center}

\paragraph{Turbulence convection is first order on purpose.} Second-order
$k$/$\omega$ transport is the usual reason residuals stall on a mesh of this
kind, and the effect on the integrated forces is small. Momentum convection
is second order.

\paragraph{Differences against the previous HELYX setup.} Relevant whenever
the two are compared:
\begin{itemize}
\item \bad{No curvature correction} (HELYX used Menter--Smirnov). The
  endplate and wheel vortex cores will diffuse somewhat faster.
\item \bad{No separation fix} (HELYX used the Rumsey form). Separation onset
  on the suction side may differ.
\item The adaptive wall treatment \emph{is} reproduced, by the all-$y^+$ wall
  functions above.
\item HELYX ran a coupled solver pseudo-transiently; this is segregated and
  steady, so iteration counts are not comparable.
\end{itemize}

% ===========================================================================
\section{Convergence}

\figOrNote{conv_residuals}{0.95}{Initial residuals per SIMPLE iteration, with
the \texttt{residualControl} thresholds marked.}{convres}

\figOrNote{conv_forces}{0.95}{Force history. The grey band is the window
averaged for the reported result.}{convf}

\figOrNote{conv_forces_zoom}{0.95}{Settling of the total $C_L\!\cdot\!A$ over
the final iterations.}{convfz}

\paragraph{How to judge this.}
\begin{itemize}
\item \texttt{residualControl} satisfied: \textbf{\converged}.
\item Final initial residuals: $p$ \resp, $U_x$ \resUx, $U_z$ \resUz,
  $k$ \resk, $\omega$ \resomega.
\item \textbf{The force spread matters more than the residual.} A converged
  residual with $C_L\!\cdot\!A$ still swinging by several percent means an
  unsteady flow being forced into a steady solve. Here the total spread over
  the averaging window is \spreadTotal\,\%.
\item Residual \emph{fields} are written as well, so a stalled case can be
  opened in ParaView to see \emph{where} it will not settle.
\end{itemize}

% ===========================================================================
\section{Results}\label{sec:results}

\paragraph{Colour ranges are fixed, not auto-scaled.} The raw fields contain
isolated extremes far outside the physical range --- measured on this build,
$C_p$ down to $-32$, $C_pT$ up to $100$ and $|U|$ up to \SI{100}{m/s} against
a \SI{10}{m/s} freestream. Those are single cells, not regions: the velocity
figure below shows nothing at the top of its scale. Auto-scaling to them
would flatten every plot, so the ranges are the physically meaningful ones
and anything outside saturates.

\figOrNote{res_umag}{0.95}{$|U|$ on the wheel-centre plane. Acceleration over
the wheel crown, the separated wheel wake, and the wing wake below. Nothing
reaches the top of the scale, which is how you tell the extreme values are
isolated cells rather than a region.}{resumag}

\figOrNote{res_cp_wing}{0.92}{$C_p$ seen from below: wing suction side and
the underside of the wheel.}{rescp}

\figOrNote{res_cp_wing_top}{0.92}{$C_p$ seen from above: wing pressure side
and the wheel stagnation region.}{rescptop}

\figOrNote{res_cpt_slice}{0.95}{Total pressure coefficient on streamwise
slices. $C_pT\approx1$ in the freestream; the dark regions are total-pressure
loss, and the discrete cores downstream are the endplate and wheel
vortices.}{rescpt}

\figOrNote{res_q}{0.92}{Q-criterion isosurface, coloured by streamwise
vorticity --- the endplate and wheel vortex systems.}{resq}

\figOrNote{res_tau}{0.92}{Wall shear stress magnitude on the suction side.
Dark bands are separation lines.}{restau}

\figOrNote{res_yplus}{0.92}{$y^+$ on the wall patches.}{resyplus}

\paragraph{$y^+$, the check that decides whether the wall treatment holds.}
Mean values: wing \yplusWing, wheel \yplusWheel. Per-patch averages on this
build sit between 2.3 and 4 on the wing and between 2.9 and 3.7 on the wheel
tread, shoulders and sidewall --- the range the layer stack was sized for.
Two outliers, both expected:
\begin{itemize}
\item \texttt{wheel-plinth} averages 8.7. It carries no prism layers, so the
  first cell is a bulk cell.
\item \texttt{ground} peaks at \yplusMax, in the contact-patch region where
  the belt passes under the wheel. Local, and not where the forces come from.
\end{itemize}

% ===========================================================================
\section{Finishing this report}\label{sec:finish}

The solve in this build stopped at \nIter\ of \nIterations\ iterations. To
produce the final version:

\begin{verbatim}
cd RANS_Simulations/S1.42_h0.13_AOA8_W0.63
./Allsolve --continue          # resumes from the latest written time
./Allpost --summary            # refreshes results.json

cd ../../report
./build.sh                     # regenerates every figure, rebuilds the PDF
\end{verbatim}

\begin{itemize}
\item \texttt{--continue} resumes rather than restarting:
  \texttt{startFrom} is \texttt{latestTime} in \texttt{controlDict}.
\item The run stops by itself once \texttt{residualControl} is met, which is
  usually well before the configured \nIterations.
\item To stop it cleanly and keep the state resumable, set
  \texttt{stopAt writeNow} in \texttt{system/controlDict} --- the solver
  writes the current state and exits. \texttt{runTimeModifiable} is on, so it
  picks the change up within one iteration. Killing the process instead
  loses everything since the last write.
\item \texttt{build.sh --no-figures} rebuilds only the text, which is quick.
\end{itemize}

\noindent
Nothing in the report is typed by hand: all numbers come from the case's
\texttt{results.json} and solver logs via
\texttt{report/scripts/make\_tex\_data.py}, so a rebuild picks up the final
solution automatically.

% ===========================================================================
\section{Assessment}

\paragraph{What supports trusting these numbers.}
\begin{itemize}
\item \ok{Background mesh exactly as designed} (\nBaseCells\ cells), cubic,
  dividing the domain without remainder.
\item \ok{Refinement reproduces the previous validated HELYX setup}
  level for level, and the final cell count agrees with it closely.
\item \ok{Patch groups resolve correctly} ---
  \texttt{patchSummary} confirms \texttt{noSlip} on the wing,
  \texttt{rotatingWallVelocity} on the wheel and the belt velocity on the
  plinth.
\item \ok{Wheel rotation sign verified} analytically against the belt speed.
\item \ok{Divergence-free initialisation} via \texttt{potentialFoam}.
\item \ok{$y^+$ in the design range} of the chosen wall functions.
\end{itemize}

\paragraph{What to be careful about.}
\begin{itemize}
\item \bad{No curvature correction and no separation fix} relative to
  HELYX. Expect faster vortex diffusion and possibly different separation
  onset.
\item \bad{Relaxed mesh quality limits} leave genuinely poor cells in the
  mesh (\texttt{minTetQuality} disabled). Accepted in exchange for layer
  coverage.
\item \bad{Three patches without prism layers}, see section~5.
\item \bad{Two duplicate triangles} in the wing STL.
\item \bad{Steady RANS on a flow that may not be steady.} At high angle of
  attack or low ride height, check the force spread before quoting a number.
\item The two local refinement boxes that stand in for HELYX's per-region
  distance refinement are \bad{positioned for this geometry} and do not
  follow the parameters. Re-check them before extending the sweep.
\end{itemize}

\vfill
\noindent\footnotesize
Generated from \texttt{\caseName}. Figures by \texttt{pvbatch}
(ParaView) and \texttt{matplotlib}; numbers by
\texttt{report/scripts/make\_tex\_data.py} from the case's
\texttt{results.json} and solver logs.

\end{document}
